\chapter{背景と目的}

\section{タスクの目的}

廃プラスチックボトルの再資源化では、洗浄、材質別分別、再利用の前に、ボトル本体とキャップを分離する必要がある。本プロジェクトでは、机上の個数、位置、向きが一定でないボトルに対して、「ボトル把持―開栓―分別投入」を連続して行うシステムを構築する。単発の展示動作ではなく、データ追加、再学習、長時間運用、定量評価を継続できることを目標とする。

\plain{人はボトルを見れば、キャップの有無、口部の方向、滑りにくい把持位置を自然に判断する。ロボットは画像と過去の示教から同じ判断を学び、二つの腕で異なる動作を同時に協調させる必要がある。}

\section{ALOHA を採用した理由}

ALOHA の中核的な利点は、双腕作業と双腕示教を同一設備で行える点にある。操作者は前側の二本の示教腕を用いて自然に動作を提示し、後側の二本の作業腕が同期して動く。同時に、多視点画像と関節動作を記録できる。初期の ALOHA は精密な双腕操作の収集に適することを示し、ALOHA 2 は操作性、堅牢性、大規模収集能力を改善した\cite{zhao2023aloha,aloha2_2024}。

本タスクに適する理由は次の三点である。

\begin{enumerate}
  \item 一方の作業腕でボトル本体を保持し、他方でキャップを把持・回転できる。
  \item 上方視点で机全体を確認し、手首視点で腕に隠れやすいボトル口とグリッパ周辺を確認できる。
  \item 連続回転、持ち直し、分別投入を直接示教でき、実データを規模化しやすい。
\end{enumerate}

\figfull[.96\textwidth]{aloha_formal_photo_annotated.png}{実機設備の写真。撮影時のロボットは静止・休止状態であり、タスク実行例や成功例ではない。前側二本が操作者用示教腕、後側二本がロボット作業腕である。}

\section{難しさは「一回つかむ」ことにとどまらない}

\figfull[.98\textwidth]{task_timeline.pdf}{ボトル 1 本を処理する 7 段階の双腕タスク。探索、把持、キャップ位置決め、回転、分離、投入、復帰は相互に依存し、前段の誤差が後段へ伝播する。現段階の達成状況は現場観察による。}

主な難しさは次の結合にある。

\begin{itemize}
  \item \textbf{対象が固定されない：}ボトルは密集・遮蔽し、向きも変化し、キャップがない場合もある。
  \item \textbf{双腕同期が必要：}左腕の保持が弱いと本体が共回りし、保持位置が悪いと口部を遮る。
  \item \textbf{接触が不確実：}キャップごとに直径、表面、摩擦、締まり具合が異なる。
  \item \textbf{多段階タスクで誤差が蓄積：}把持誤差が位置決め、回転、分離、投入へ影響する。
  \item \textbf{長時間循環が必要：}一回の成功は継続運転を意味せず、復帰誤差と異常復旧が時間とともに蓄積する。
\end{itemize}

\section{本報告書の範囲}

本報告書は、目的、方法、データ、モデル、実験、結果、次年度計画を順に整理する。結論には、学習履歴、学習済みモデル、データ統計、設備写真、示教フレーム、アテンション重みの記録、現場観測のみを用いる。動画または計数表がない内容は現場概算として扱う。3 年間の全体目標は扱わず、将来の水管挿入洗浄を完了済み能力として記載しない。
